\chapter{双线性函数与辛空间}
\section{习题}
\begin{enumerate}
  \item 设 $V$ 是数域 $P$ 上一个 3 维线性空间，$\varepsilon_1,\varepsilon_2,\varepsilon_3$ 是它的一组基，$f$ 是 $V$ 上一个线性函数，已知
  \[
    f(\varepsilon_1+\varepsilon_3)=1,\quad f(\varepsilon_2-2\varepsilon_3)=-1,\quad f(\varepsilon_1+\varepsilon_2)=-3,
  \]
  求 $f(x_1\varepsilon_1+x_2\varepsilon_2+x_3\varepsilon_3)$．

  \anscompute
  \item $V$ 及 $\varepsilon_1,\varepsilon_2,\varepsilon_3$ 同上题，试找出一个线性函数 $f$，使
  \[
    f(\varepsilon_1+\varepsilon_3)=f(\varepsilon_1-2\varepsilon_3)=0,\quad f(\varepsilon_1+\varepsilon_2)=1.
  \]

  \ansgeneral
  \item 设 $\varepsilon_1,\varepsilon_2,\varepsilon_3$ 是线性空间 $V$ 的一组基，$f_1,f_2,f_3$ 是它的对偶基，
  \[
    \alpha_1=\varepsilon_1-\varepsilon_3,\quad \alpha_2=\varepsilon_1+\varepsilon_2+\varepsilon_3,\quad \alpha_3=\varepsilon_2+\varepsilon_3.
  \]
  试证 $\alpha_1,\alpha_2,\alpha_3$ 是 $V$ 的一组基并求它的对偶基（用 $f_1,f_2,f_3$ 表出）．

  \ansproof
  \item 设 $V$ 是一个线性空间，$f_1,f_2,\cdots,f_s$ 是 $V^*$ 中非零向量．试证：存在 $\alpha\in V$，使
  \[
    f_i(\alpha)\neq 0,\quad i=1,2,\cdots,s.
  \]

  \ansproof
  \item 设 $\alpha_1,\alpha_2,\cdots,\alpha_s$ 是线性空间 $V$ 中非零向量．证明：存在 $f\in V^*$，使
  \[
    f(\alpha_i)\neq 0,\quad i=1,2,\cdots,s.
  \]

  \ansproof
  \item 设 $V=P[x]_3$，对 $p(x)=c_0+c_1x+c_2x^2\in V$ 定义
  \[
    f_1(p(x))=\int_0^1 p(x)\,\mathrm{d}x,\quad
    f_2(p(x))=\int_0^2 p(x)\,\mathrm{d}x,\quad
    f_3(p(x))=\int_0^{-1} p(x)\,\mathrm{d}x.
  \]
  试证 $f_1,f_2,f_3$ 都是 $V$ 上线性函数，并找出 $V$ 的一组基 $p_1(x),p_2(x),p_3(x)$，使 $f_1,f_2,f_3$ 是它的对偶基．

  \ansproof
  \item 设 $V$ 是一个 $n$ 维欧氏空间，它的内积为 $(\alpha,\beta)$，对 $V$ 中确定的向量 $\alpha$，定义 $V$ 上一个函数 $\alpha^*$：
  \[
    \alpha^*(\beta)=(\alpha,\beta).
  \]
  \begin{enumerate}
    \item 证明：$\alpha^*$ 是 $V$ 上线性函数；
  \anssub
    \item 证明：$V$ 到 $V^*$ 的映射 $\alpha\to\alpha^*$ 是 $V$ 到 $V^*$ 的一个同构映射．（在这个同构下，欧氏空间可看成自身的对偶空间．）
  \anssub
  \end{enumerate}

  \item 设 $\mathcal{A}$ 是 $P$ 上 $n$ 维线性空间 $V$ 的一个线性变换．
  \begin{enumerate}
    \item 证明：对 $V$ 上的线性函数 $f$，$f\mathcal{A}$ 仍是 $V$ 上线性函数；
  \anssub
    \item 定义 $V^*$ 到自身的映射 $\mathcal{A}^*$ 为
    \[
      f\to f\mathcal{A},
    \]
    证明：$\mathcal{A}^*$ 是 $V^*$ 上的线性变换；
  \anssub
    \item 设 $\varepsilon_1,\varepsilon_2,\cdots,\varepsilon_n$ 是 $V$ 的一组基，$f_1,f_2,\cdots,f_n$ 是它的对偶基，并设 $\mathcal{A}$ 在 $\varepsilon_1,\varepsilon_2,\cdots,\varepsilon_n$ 下的矩阵为 $A$．证明：$\mathcal{A}^*$ 在 $f_1,f_2,\cdots,f_n$ 下的矩阵为 $A^{\mathrm{T}}$．（因此 $\mathcal{A}^*$ 称为 $\mathcal{A}$ 的转置映射．）
  \anssub
  \end{enumerate}

  \item 设 $V$ 是数域 $P$ 上一个线性空间，$f_1,f_2,\cdots,f_k$ 是 $V$ 上 $k$ 个线性函数．
  \begin{enumerate}
    \item 证明集合
    \[
      W=\{\alpha\in V\mid f_i(\alpha)=0,\ 1\leq i\leq k\}
    \]
    是 $V$ 的一个子空间，$W$ 称为线性函数 $f_1,f_2,\cdots,f_k$ 的零化子空间；
  \anssub
    \item 证明：$V$ 的任意一个子空间皆为某些线性函数的零化子空间．
  \anssub
  \end{enumerate}

  \item 设 $A$ 是 $P$ 上一个 $m$ 阶矩阵，定义 $P^{m\times n}$ 上一个二元函数
  \[
    f(X,Y)=\operatorname{tr}(X^{\mathrm{T}}AY)=X^{\mathrm{T}}AY\text{ 的主对角线上元素的和},\quad X,Y\in P^{m\times n}.
  \]
  \begin{enumerate}
    \item 证明：$f(X,Y)$ 是 $P^{m\times n}$ 上的双线性函数；
  \anssub
    \item 求 $f(X,Y)$ 在基
    \[
      E_{11},E_{12},\cdots,E_{1n},E_{21},E_{22},\cdots,E_{2n},\cdots,E_{m1},E_{m2},\cdots,E_{mn}
    \]
    下的度量矩阵．（$E_{ij}$ 表示第 $i$ 行第 $j$ 列的元素为 $1$，而其余元素全为零的 $m\times n$ 矩阵．）
  \anssub
  \end{enumerate}

  \item 在 $P^4$ 中定义一个双线性函数 $f(X,Y)$，对 $X=(x_1,x_2,x_3,x_4)$，$Y=(y_1,y_2,y_3,y_4)$，$f(X,Y)=3x_1y_2-5x_2y_1+x_3y_4-4x_4y_3$．
  \begin{enumerate}
    \item 给定 $P^4$ 的一组基
    \[
      \begin{aligned}
        \varepsilon_1&=(1,-2,-1,0),&\quad \varepsilon_2&=(1,-1,1,0),\\
        \varepsilon_3&=(-1,2,1,1),&\quad \varepsilon_4&=(-1,-1,0,1).
      \end{aligned}
    \]
    求 $f(X,Y)$ 在这组基下的度量矩阵；
  \anssub
    \item 另取一组基 $\eta_1,\eta_2,\eta_3,\eta_4$，且
    \[
      (\eta_1,\eta_2,\eta_3,\eta_4)=(\varepsilon_1,\varepsilon_2,\varepsilon_3,\varepsilon_4)T,
    \]
    其中
    \[
      T=\begin{pmatrix}
        1 & 1 & 1 & 1 \\
        1 & 1 & -1 & -1 \\
        1 & -1 & 1 & -1 \\
        1 & -1 & -1 & 1
      \end{pmatrix},
    \]
    求 $f(X,Y)$ 在 $\eta_1,\eta_2,\eta_3,\eta_4$ 下的度量矩阵．
  \anssub
  \end{enumerate}

  \item 设 $V$ 是复数域上线性空间，其维数 $n\geq 2$，$f(\alpha,\beta)$ 是 $V$ 上一个对称双线性函数．
  \begin{enumerate}
    \item 证明：$V$ 中有非零向量 $\xi$，使
    \[
      f(\xi,\xi)=0;
    \]
  \anssub
    \item 如果 $f(\alpha,\beta)$ 是非退化的，则必有线性无关的向量 $\xi,\eta$ 满足
    \[
      f(\xi,\eta)=1,\quad f(\xi,\xi)=f(\eta,\eta)=0.
    \]
  \anssub
  \end{enumerate}

  \item 试证线性空间 $V$ 上双线性函数 $f(\alpha,\beta)$ 为反称的充要条件是对任意 $\alpha\in V$，都有 $f(\alpha,\alpha)=0$．

  \ansproof
  \item 设 $f(\alpha,\beta)$ 是 $V$ 上对称的或反称的双线性函数，$\alpha,\beta$ 是 $V$ 中两个向量，\\
  如果 $f(\alpha,\beta)=0$，则称 $\alpha,\beta$ 正交．\par
  再设 $K$ 是 $V$ 的一个真子空间．证明：对 $\xi\notin K$，必有 $0\neq\eta\in K+L(\xi)$，使
  \[
    f(\eta,\alpha)=0
  \]
  对所有 $\alpha\in K$ 都成立．

  \ansproof
  \item 设 $V$ 与 $f(\alpha,\beta)$ 同上题，$K$ 是 $V$ 的一个子空间，令
  \[
    K^{\perp}=\{\alpha\in V\mid f(\alpha,\beta)=0,\ \forall\beta\in K\}.
  \]
  \begin{enumerate}
    \item 试证 $K^{\perp}$ 是 $V$ 的子空间（$K^{\perp}$ 称为 $K$ 的正交补）；
  \anssub
    \item 试证，如果 $K\cap K^{\perp}=\{0\}$，则 $V=K+K^{\perp}$．
  \anssub
  \end{enumerate}

  \item 设 $V,f(\alpha,\beta),K$ 同上题，并设 $f(\alpha,\beta)$ 限制在 $K$ 上是非退化的．试证：$V=K+K^{\perp}$，并证明 $f(\alpha,\beta)$ 在 $K^{\perp}$ 上是非退化的充要条件是 $f(\alpha,\beta)$ 在 $V$ 上为非退化的．

  \ansproof
  \item 设 $f(\alpha,\beta)$ 是 $n$ 维线性空间 $V$ 上的非退化对称双线性函数，对 $V$ 中一个元素 $\alpha$，定义 $V^*$ 中一个元素 $\alpha^*$：
  \[
    \alpha^*(\beta)=f(\alpha,\beta),\quad \beta\in V.
  \]
  试证：
  \begin{enumerate}
    \item $V$ 到 $V^*$ 的映射
    \[
      \alpha\to\alpha^*
    \]
    是一个同构映射；
  \anssub
    \item 对 $V$ 的每组基 $\varepsilon_1,\varepsilon_2,\cdots,\varepsilon_n$，有 $V$ 的唯一的一组基 $\varepsilon_1',\varepsilon_2',\cdots,\varepsilon_n'$，使
    \[
      f(\varepsilon_i,\varepsilon_j')=\delta_{ij};
    \]
  \anssub
    \item 如果 $V$ 是复数域上 $n$ 维线性空间，则有一组基 $\eta_1,\eta_2,\cdots,\eta_n$，使
    \[
      \eta_i=\eta_i',\quad i=1,2,\cdots,n.
    \]
  \anssub
  \end{enumerate}

  \item 设 $V$ 是对于非退化对称双线性函数 $f(\alpha,\beta)$ 的 $n$ 维准欧氏空间．如果 $V$ 的一组基 $\varepsilon_1,\varepsilon_2,\cdots,\varepsilon_n$ 满足
  \[
    \begin{cases}
      f(\varepsilon_i,\varepsilon_i)=1, & i=1,2,\cdots,p;\\
      f(\varepsilon_i,\varepsilon_i)=-1, & i=p+1,\cdots,n;\\
      f(\varepsilon_i,\varepsilon_j)=0, & i\neq j,
    \end{cases}
  \]
  则称为 $V$ 的一组正交基．如果 $V$ 上的线性变换 $\mathcal{A}$ 满足
  \[
    f(\mathcal{A}\alpha,\mathcal{A}\beta)=f(\alpha,\beta),\quad \alpha,\beta\in V,
  \]
  则称 $\mathcal{A}$ 为 $V$ 的一个准正交变换．试证：
  \begin{enumerate}
    \item 准正交变换是可逆的，且逆变换也是准正交变换；
  \anssub
    \item 准正交变换的乘积仍是准正交变换；
  \anssub
    \item 准正交变换 $\mathcal{A}$ 的特征向量 $\alpha$ 若满足 $f(\alpha,\alpha)\neq 0$，则其特征值等于 $1$ 或 $-1$；
  \anssub
    \item 准正交变换在正交基下的矩阵 $T$ 满足
    \[
      T^{\mathrm{T}}
      \begin{pmatrix}
        1 & & & & & \\
         & \ddots & & & & \\
         & & 1 & & & \\
         & & & -1 & & \\
         & & & & \ddots & \\
         & & & & & -1
      \end{pmatrix}
      T=
      \begin{pmatrix}
        1 & & & & & \\
         & \ddots & & & & \\
         & & 1 & & & \\
         & & & -1 & & \\
         & & & & \ddots & \\
         & & & & & -1
      \end{pmatrix}.
    \]
  \anssub
  \end{enumerate}
\end{enumerate}
